
% --- Title block ---
\begin{center}
  {\LARGE\bfseries rvLLM: Single-GPU, FP8, Graph-Captured\par LLM Inference in Rust\par}
  \vspace{0.6em}
  {\large Andy Norris (m0at)\par}
  \vspace{0.3em}
  {\normalsize\texttt{https://github.com/m0at/rvllm}\par}
  \vspace{0.3em}
  {\normalsize San Francisco, April 16, 2026\par}
\end{center}
\vspace{1em}

% --- Abstract ---
\begin{abstract}
rvLLM is a single-GPU LLM inference engine written in Rust, built from
scratch without Python in the serving hot path. On a single NVIDIA H100 SXM
80\,GB running Qwen2.5-7B in FP8~E4M3 with 512 output tokens, greedy decoding,
and CUDA graphs enabled, rvLLM reaches \textbf{19{,}287 tokens/second at
$N\!=\!128$} (SHA \texttt{eb9e247fd}, 3 iterations, stdev~38.5).

The engine ships as a FlashAttention-3 SM90 paged-KV decode kernel (built from
the Hopper source to \texttt{libfa3\_kernels.so}), a CUTLASS 3.x SM90 FP8 GEMM
variant catalog (40~non-residual + 10~residual-fused templates), a set of
eight custom fused PTX kernels (norm+quantize, residual+norm+quantize,
RoPE+KV-write, SiLU*mul+quantize, argmax, embedding, residual-add,
post-attention quantize), and a graph-capture runtime that records
\textbf{311 kernel launches per decode step into a single \texttt{cuGraphLaunch}}
for each of 35 pre-captured batch buckets. Total device-to-host traffic per
decode step is one $4\!\cdot\!N$-byte copy of argmax'd token identifiers;
total host-to-device traffic is one packed metadata buffer
(${<}100$\,KB at $N\!=\!128$).

The paper states the engine's five correctness invariants and enumerates
the three structural bug classes that the engine is designed to eliminate:
schedule-epilogue pairing mismatches in SM90 CUTLASS FP8 GEMMs (silent
\texttt{CUDA\_ERROR\_ILLEGAL\_ADDRESS} under graph replay only), captured
metadata offsets clobbered by orthogonal upload paths, and stale paged-KV
block tables from CoW-only change detection. Each is eliminated by a
pinned variant, a frozen layout, or a typed signal from the scheduler.
Missing build artifacts (autotune policy, FA3~\texttt{.so},
CUTLASS~\texttt{.so}) refuse to start; no silent fallbacks exist in the
FP8 dispatch. rvLLM is open-source under Apache-2.0.
\end{abstract}

% ============================================================================
\section{Introduction}
\label{sec:introduction}
% ============================================================================

Inference serving for large language models is bandwidth-bound at decode
batch sizes of practical interest. On an H100 running Qwen2.5-7B, a single
decode step involves roughly 10--15 kernel launches per transformer layer,
paged-KV attention over the sequence's current context, and a final
projection to a 152{,}064-entry vocabulary. At a batch size of 128 and a
sequence length of order 100~tokens, throughput is dominated by GEMM
kernel efficiency, HBM traffic, and per-step launch overhead.

vLLM~\cite{kwon2023vllm} established the paged-KV cache design and
continuous batching scheduler that define the modern serving baseline.
CUTLASS~3.x~\cite{cutlass2024} and FlashAttention-3~\cite{shah2024flash3}
supplied the kernel primitives that make Hopper-specific WGMMA and TMA
accessible to non-Python callers. rvLLM composes these primitives behind
a Rust runtime that owns CUDA graph capture, paged KV allocation,
metadata packing, and sampling DtoH coordination in 31 crates
($\sim$76K lines of Rust), with no Python interpreter in the serving
path.

This paper documents the state of the engine at commit
\texttt{eb9e247fd} (April~16, 2026), which achieves
19{,}287~tok/s at $N\!=\!128$. Section~\ref{sec:stack} walks the full
stack from HTTP handler to kernel. Section~\ref{sec:kernels} enumerates
the kernel inventory. Section~\ref{sec:step} gives an ordered trace of
one decode step. Section~\ref{sec:invariants} states the five correctness
invariants the engine enforces and ties each to a concrete failure mode
observed during development. Section~\ref{sec:recovery} describes three
root-cause commits that, on April~16, took the engine from
9{,}531~tok/s to 19{,}287~tok/s, none of which were stopgaps.
Section~\ref{sec:results} reports the full bench table; related work and
limitations follow.

% ============================================================================
\section{The stack, every layer}
\label{sec:stack}
% ============================================================================

\begin{figure}[h]
\footnotesize
\begin{verbatim}
HTTP / OpenAI API              rvllm-api (axum)
Engine                         rvllm-v2::engine
                               step_pipelined() = launch + collect
                               double-buffered pinned argmax DtoH
Scheduler                      continuous batching
                               paged KV (block_size = 64)
                               page-boundary growth signal
Worker                         CUDA graph pool
                               35 pre-captured batch buckets
                               1 compute stream per worker
Runner                         1 HBM arena, pre-allocated slab
                               packed metadata (1 HtoD per step)
Layer                          11 kernel launches per layer
Kernels                        CUTLASS 3.x SM90 FP8 GEMM
                               FlashAttention-3 SM90 paged decode
                               custom fused PTX (norm / silu / rope)
                               cuBLAS HGEMM (LM head)
\end{verbatim}
\caption{The rvLLM stack. Each row runs in a single Rust crate; the DAG
is enforced by a compile-time test that parses every \texttt{Cargo.toml}
and rejects out-of-layer dependencies.}
\label{fig:stack}
\end{figure}

\subsection{HTTP, engine, scheduler}

\textbf{HTTP} is an \texttt{axum} server exposing OpenAI-compatible
\texttt{/v1/completions}, \texttt{/v1/chat/completions},
\texttt{/v1/embeddings}, and \texttt{/v1/responses}. Incoming requests are
tokenized and placed on a lock-free MPSC queue consumed by a single worker
thread; the worker is \texttt{!Send} by construction so a tokio runtime
cannot accidentally schedule two concurrent GPU users.

\textbf{The engine} (\texttt{rvllm-v2::engine}) exposes two functions:
\texttt{step\_launch(diff)} enqueues all GPU work for one step (metadata
HtoD, graph replay, argmax DtoH, event record) and returns immediately;
\texttt{step\_collect()} waits on the previous step's DtoH event and
returns the new tokens. There is exactly one codepath; the synchronous
variant that tied \texttt{cuStreamSynchronize} to each step has been
removed.

\textbf{The scheduler} (\texttt{rvllm-scheduler}) implements continuous
batching and paged-KV allocation. A request moves through
\texttt{Queued} $\to$ \texttt{Prefilling} $\to$ \texttt{Decoding} $\to$
\texttt{Finished}. The scheduler emits one \texttt{BatchPlan} per step,
of exactly one variant (\texttt{Prefill}, \texttt{Decode}, or
\texttt{Idle}); prefill and decode never coexist in the same step. A
watermark-triggered preemption path re-queues the lowest-priority
decoding requests when KV-cache block use exceeds 96\%.

\subsection{Worker, runner, layer}

\textbf{The worker} (\texttt{rvllm-worker} / \texttt{rvllm-v2::worker})
owns one CUDA compute stream, one CUDA context, and a pool of pre-captured
graph instances indexed by bucket size. Buckets are
$\{1,2,4,8,16,24,32,\ldots,256\}$. Graph capture happens during engine
initialization, not during warmup.

\textbf{The runner} (\texttt{rvllm-v2::runner}) owns one HBM arena
allocated via one \texttt{cuMemAlloc} at startup, sized for the
worst-case workspace (sweep over all CUTLASS variants $\times$ all
bucket sizes $\times$ all GEMM shapes), plus the KV cache, scratch
tensors, and packed metadata buffer. A compile-time \texttt{GraphSafe}
trait marks tensors whose lifetime outlives the capture region;
\texttt{\&mut HbmArena} is not in scope inside a \texttt{CaptureScope}
closure, so runtime realloc is unrepresentable.

\textbf{The layer} (\texttt{rvllm-v2::layer}) is a pure function of
(input, weights, KV cache slab, scratch, metadata handle). One decode
layer at FP8 is eleven kernel launches:

\begin{verbatim}
1.  fused_add_rmsnorm_fp8_quant
2.  CUTLASS FP8 GEMM (QKV)
3.  fused_rope_kv_write
4.  FA3 paged_decode
5.  quantize_fp8_per_token (post-attention)
6.  CUTLASS FP8 GEMM + residual (o_proj, v1)
7.  fused_rmsnorm_fp8_quant
8.  CUTLASS FP8 GEMM (gate_up)
9.  fused_silu_mul_fp8_quant
10. CUTLASS FP8 GEMM (down_proj)
11. residual_add_f16
\end{verbatim}

Buffer aliasing is checked by the borrow checker: the layer takes
\texttt{\&mut Tensor<'a, f16>} for its output and \texttt{\&mut LayerScratch}
for scratch; two mutable borrows into one allocation fail to compile.

% ============================================================================
\section{Kernel inventory}
\label{sec:kernels}
% ============================================================================

Every kernel has a pinned variant and a workspace contract. There are
no dispatch fallback chains: each call site resolves to exactly one
variant, determined at engine init from the autotune policy or pinned
in source.

\subsection{CUTLASS SM90 FP8 GEMM}

CUTLASS~3.x templates drive the QKV, gate\_up, down\_proj, and fused
o\_proj+residual projections. Forty non-residual and ten residual-fused
variants are instantiated, each parameterized by
$(\text{TileShape}, \text{ClusterShape}, \text{KernelSchedule})$.

Variants cover tile shapes
$\{64\!\times\!128, 128\!\times\!128, 128\!\times\!256\}$ in $M\!\times\!N$
with $K\!=\!128$ or $256$; cluster shapes $1\!\times\!1\!\times\!1$; and
the \texttt{KernelTmaWarpSpecialized} (WS),
\texttt{KernelTmaWarpSpecializedCooperative} (Coop), and their
\texttt{FP8FastAccum} variants. Per-row activation scale and per-tensor
weight scale are applied in the epilogue; the residual-fused kernel adds
an \texttt{Sm90SrcFetch} fetch of the residual and a \texttt{TreeVisitor}
computing $D = \alpha\cdot\text{GEMM}(A,B) + C$.

\textbf{Pairing constraint.} The mainloop schedule must match the
epilogue schedule. The residual-fused kernel uses
\texttt{TmaWarpSpecializedCooperative} in its epilogue collective
(\texttt{kernels/cutlass\_fp8\_gemm\_residual.cu}, line 100); any variant
passing a non-cooperative mainloop with this epilogue causes
\texttt{CUDA\_ERROR\_ILLEGAL\_ADDRESS} \emph{only under graph replay}
(the kernel runs cleanly outside a captured graph). The v2 runtime pins
the o\_proj residual path to variant \texttt{v1}
($128\!\times\!128\!\times\!128$ Coop/Coop matched); the v3 rewrite makes
the mismatch a CUDA \texttt{static\_assert}, rejecting the offending
template pair at compile time.

\textbf{Workspace contract.} Each variant exposes
\texttt{fp8\_gemm\_v<i>\_workspace\_size(M,N,K)}; the runtime's
\texttt{max\_workspace\_size} query iterates all actually-callable
variants across all bucket sizes for all GEMM shapes the model uses,
and pre-allocates a single workspace slab accordingly. This eliminates
the class of bug where the dispatch falls through to a variant whose
workspace was never queried.

\textbf{Autotune policy.} A separate binary, \texttt{autotune-cutlass},
runs 67~HGEMM variants, 31~o-proj-residual variants, 32~gate-up-silu
variants, 40~FP8 GEMM variants, and 10~FP8-GEMM-residual variants
against every GEMM shape the model uses at every bucket size, reporting
median of 10 timed runs after 3 warmups. The winning
$(\text{shape} \to \text{variant index})$ mapping is written to
\texttt{policy.json}, which is a build artifact shipped alongside the
binaries. The runtime does not consult \texttt{\textasciitilde/.cache}
for policy; stale caches from prior deploys cannot influence dispatch.

\subsection{FlashAttention-3 SM90 paged decode}

\texttt{libfa3\_kernels.so} is built at deploy time from the
FlashAttention-3 Hopper source tree plus an
\texttt{extern "C"} wrapper (\texttt{kernels/fa3\_sm90\_wrapper.cu}).
Build time is roughly 10~minutes on an H100; the resulting shared object
is 23~MB. The wrapper exposes \texttt{paged\_decode} and
\texttt{paged\_prefill} entry points. Both expect KV laid out as
$[2, \text{num\_blocks}, \text{block\_size}, \text{num\_kv\_heads},
\text{head\_dim}]$ in one contiguous \texttt{Region}; the factor-of-two
comes from interleaving K and V per block.

GQA is supported via $\text{num\_heads} / \text{num\_kv\_heads}$;
\texttt{head\_dim = 128} is a hard gate (the kernel is specialized and
refuses other head dimensions). \texttt{context\_lens[i] == 0} is a
valid padded-slot marker and yields zero output without touching
\texttt{block\_tables[i,*]}.

\textbf{No PTX fallback.} A prior PTX-only attention kernel existed
in the tree. It has been removed from the dispatch; if
\texttt{libfa3\_kernels.so} is not present at engine init, the engine
refuses to load. The build script ships with the repository.

\subsection{Fused pre/post kernels}

Eight custom CUDA kernels compile to PTX and are loaded by the runtime.
Each fuses one recognizable composite. No megakernels exist in the
dispatch path.

\begin{tabular}{lll}
\toprule
\textbf{Kernel} & \textbf{Fuses} & \textbf{Launches} \\
\midrule
\texttt{embedding\_gather}             & token\_ids $\to$ f16 hidden & $1\to 1$ \\
\texttt{fused\_add\_rmsnorm\_fp8\_quant} & residual + RMSNorm + quant & $3\to 1$ \\
\texttt{fused\_rmsnorm\_fp8\_quant}      & RMSNorm + quant            & $2\to 1$ \\
\texttt{quantize\_fp8\_per\_token}       & post-attention act $\to$ FP8 & $1\to 1$ \\
\texttt{fused\_rope\_kv\_write}         & RoPE + paged KV write      & $3\to 1$ \\
\texttt{fused\_silu\_mul\_fp8\_quant}    & SiLU + mul + quant         & $3\to 1$ \\
\texttt{argmax}                         & f32 logits $\to$ i32 token & $1\to 1$ \\
\texttt{residual\_add\_f16}             & $x + y$ in place           & $1\to 1$ \\
\bottomrule
\end{tabular}

Each fused kernel ships with a pure-Rust f32 reference implementation
in tests; the PTX must match within cosine~$0.999$ for FP8 outputs and
$10^{-5}$ absolute for f32 outputs, over a fuzz of shape and value
distributions. The \texttt{residual\_add\_f16} kernel is deliberately
\emph{not} fused into the GEMM epilogue; prior attempts to fuse residual
with the downstream projection hit either schedule-pairing crashes or a
$\sim$24\% regression from increased tile pressure.

\textbf{Vectorization rule.} Kernels that accept dimensions divisible by
eight halves issue \texttt{uint4} 128-bit loads (eight halves per load)
and \texttt{uint2} 64-bit FP8 stores (eight FP8 per store), with register
caching between the absolute-max reduction pass and the quantize pass so
each HBM element is read once. Alignment is checked by a
\texttt{debug\_assert} in the Rust binding; misalignment fails the launch
loudly rather than producing
\texttt{CUDA\_ERROR\_MISALIGNED\_ADDRESS} under graph replay.

% ============================================================================
\section{One decode step, in order}
\label{sec:step}
% ============================================================================

For a decode batch of $N$ sequences at bucket size
$B = \text{pad\_up}(N)$:

\begin{verbatim}
step_launch(diff):
  1. schedule(diff)                         CPU, O(N), zero allocs
  2. metadata pack + 1x HtoD                packed i32 buffer:
                                            positions, context_lens,
                                            block_tables, slot_mapping
  3. graph_pool.replay(bucket = B)          one cuGraphLaunch:
       embedding_gather                      1 kernel
       for layer in 0..28:                   11 kernels x 28 = 308
         (as enumerated in Section 2.2)
       fused_residual_rmsnorm (final)        1 kernel
       cuBLAS HGEMM -> lm_head               1 kernel
       argmax                                1 kernel
  4. cuMemcpyDtoHAsync argmax -> pinned[w]  4*N bytes (512 B at N=128)
     cuEventRecord event[w]
  5. return; w ^= 1                         double-buffer flip

step_collect():
  1. cuEventSynchronize event[r]            waits for step-1 DtoH
  2. read pinned[r][0..N]                   CPU reads new tokens
  3. scheduler.commit()                     mark tokens, free finished
\end{verbatim}

Total kernel launches per decode step: $1 + 28\!\times\!11 + 1 + 1 + 1
= \mathbf{311}$, all captured into \textbf{one} \texttt{cuGraphLaunch}.
Total DtoH per step: one $4\!\cdot\!N$ byte copy of argmax token
identifiers. Total HtoD per step: one packed metadata buffer,
${<}100$\,KB at $N\!=\!128$ with $\text{max\_blocks}\!=\!129$.

The LM-head GEMM uses cuBLAS HGEMM (F16 in, F32 out) because CUTLASS
FP8 autotune reported worse time on the $(N, 152064, 3584)$ shape at
$M\!\ge\!128$ than cuBLAS for the model tested. An FP8 LM-head path
exists in \texttt{rvllm-v2} gated on the autotune cache's preferred
variant; it is disabled by default for Qwen2.5-7B.

% ============================================================================
\section{Correctness invariants}
\label{sec:invariants}
% ============================================================================

The engine enforces five invariants at build or startup time. Each is
a response to a concrete failure mode observed during development.

\subsection{No fallbacks}

\textbf{Rule.} Missing autotune-policy entry for a shape panics with the
shape. Missing \texttt{libfa3\_kernels.so} refuses startup. Missing
\texttt{libcutlass\_kernels.so} refuses startup. The runtime does not
read \texttt{\textasciitilde/.cache/rvllm/cutlass\_autotune.json}; policy
lives in the SHA-pinned \texttt{policy.json} shipped alongside the
binary.

\textbf{Failure mode this prevents.} A prior deploy populated the
on-box autotune cache with entries for an older CUTLASS catalog. A
fresh build with the same HF-downloaded kernels inadvertently consulted
the stale cache, resolving shape $(32, 152064, 3584)$ to a variant
whose template no longer matched the \texttt{.so}'s symbol table, and
ran the corresponding FP8FastAccum kernel under graph capture. Under
graph replay, the kernel faulted with
\texttt{CUDA\_ERROR\_ILLEGAL\_ADDRESS}. SHA-pinning the policy
artifact eliminates the failure class.

\subsection{Single metadata-upload path}

\textbf{Rule.} Metadata buffer layout is frozen per
$(\text{bucket}, \text{max\_blocks\_per\_seq})$. Captured graphs bind
those exact device offsets. There is no non-padded upload variant.
Prefill and decode use separate entry points that do not share offsets.

\textbf{Failure mode.} During development, two upload paths coexisted:
\texttt{upload\_metadata} (non-padded, sized to the actual request
count) and \texttt{upload\_metadata\_padded} (sized to the captured
bucket). On the first decode step after a prefill, the runtime's
\texttt{last\_meta\_offsets} field carried the non-padded offsets from
the prefill call, but the captured graph still read from the padded
offsets. The diff-upload path wrote new token IDs, positions, and
block-table entries at the non-padded indices; the captured kernel
interpreted the bytes at the padded indices as if they were current
data, producing a \texttt{context\_lens} value of order $10^5$, a
\texttt{total\_tiles} of $780$, and a \texttt{block\_tables} read 391
entries past the slice end. The read resolved to an uninitialized
region of the meta buffer, which held bytes of a small
integer interpreted as \texttt{phys\_block}. Multiplied by the
per-block stride $16384$, this yielded a device address 10.2~MB before
the nearest valid allocation---verbatim the
\texttt{compute-sanitizer} report of the original crash.

The v3 rewrite removes non-padded upload entirely. The v2 engine
tracks \texttt{last\_padded\_batch: Option<usize>} and only takes the
patch fast-path when it matches the requested bucket; otherwise a
full padded upload happens.

\subsection{Real block-change detection}

\textbf{Rule.} The scheduler emits
\texttt{ContinuedRequest::block\_table\_update:} \texttt{Option<Vec<BlockId>>}
whenever a sequence's physical block list has grown since the last
\texttt{mark\_table\_sent()}. The worker unions this with
\texttt{diff.block\_ops.copies} (copy-on-write events) to compute
\texttt{has\_block\_changes} before each patch-decode upload.

\textbf{Failure mode.} An earlier version of the worker checked only
\texttt{!diff.block\_ops.copies.is\_empty()}, which captures CoW
events but \emph{not} ordinary page-boundary growth. A sequence
reaching the end of its current block and being allocated a new
physical block generates no CoW event; the block-table entry at index
$\lceil \text{context\_len} / \text{block\_size} \rceil - 1$ changes
silently. The patch-decode path skipped the \texttt{block\_tables}
copy; the captured graph continued to read the stale block identifier,
producing wrong KV-cache reads, garbage attention outputs, and
garbage tokens. No crash. \texttt{ignore\_eos=true} bench runs at
output\_len~512 still reported plausible throughput; correctness was
degraded silently. The scheduler already emitted the needed signal;
the worker simply did not consult it.

\subsection{CUTLASS schedule/epilogue pairing}

\textbf{Rule.} Mainloop schedule and epilogue schedule must match. A
residual-fused kernel with \texttt{TmaWarpSpecializedCooperative}
epilogue requires a cooperative mainloop schedule.

\textbf{Enforcement.} The v3 catalog uses a CUDA
\texttt{static\_assert} on the paired schedule types at template
instantiation; the variant table generator writes Rust and CUDA sides
together, so adding a mismatched variant fails compilation. In v2, the
o\_proj residual dispatch pins to the hand-verified variant \texttt{v1}
($128\!\times\!128\!\times\!128$ Coop/Coop), and the catalog has been
audited for pairings---the offending variants are
$\text{cutlass\_fp8\_gemm\_v}\{0,2,3,4,6,8,11,12,13,14\}$ and
$\text{cutlass\_fp8\_gemm\_residual\_v}\{0,2,4,5,6\}$.

\subsection{No \texttt{unwrap()} or \texttt{String} errors across crates}

\textbf{Rule.} All errors crossing crate boundaries are
\texttt{Result<T, RvllmError>}. The error variants carry structured
context---stream identifier, kernel name, launch configuration, GEMM
shape, workspace size needed vs given---not a stringified
\texttt{DriverError}.

\textbf{Failure mode.} During the April~15 debugging pass, several
downstream error paths presented the root-cause \texttt{cuGraphLaunch
failed} as the opaque string
\texttt{"dtoh event sync: DriverError(CUDA\_ERROR\_ILLEGAL\_ADDRESS, \dots)"}.
The failing kernel name, stream identifier, bucket size, and
kernel-launch arguments were discarded at the first \texttt{map\_err}.
Recovery required \texttt{compute-sanitizer} to re-derive what the
error path should have carried.

% ============================================================================
\section{Measured recovery, April 16, 2026}
\label{sec:recovery}
% ============================================================================

Three root-cause fixes, in three commits, took the engine from
9{,}531~tok/s to 19{,}287~tok/s at $N\!=\!128$ in one working day. None
were stopgaps; each eliminated a structural failure mode documented in
Section~\ref{sec:invariants}.

\begin{tabular}{lll}
\toprule
\textbf{Commit} & \textbf{Fix} & $\boldsymbol{\Delta\text{(N=128)}}$ \\
\midrule
\texttt{6dabb76a5}
  & Track \texttt{last\_padded\_batch}; only patch metadata when bucket matches
  & crash $\to$ runs \\
\texttt{31716269d}
  & Re-enable fused FP8 o\_proj + residual on variant \texttt{v1} (Coop/Coop)
  & +83\% \\
\texttt{eb9e247fd}
  & Real block-change detection via \texttt{block\_table\_update} (not CoW-only)
  & +10\% \\
\bottomrule
\end{tabular}

In addition, \texttt{libfa3\_kernels.so} was built on the deploy host
(23~MB, once), eliminating a silent \texttt{.ptx} attention fallback path.
All four changes taken together restore the engine to a configuration
where every decode step uses the intended kernels, with no silent
degradation in the dispatch.

% ============================================================================
\section{Results}
\label{sec:results}
% ============================================================================

\subsection{Experimental setup}

All measurements in this section are on a single NVIDIA H100~SXM~80\,GB
(driver compatible with CUDA~12.4), running Qwen2.5-7B from the public
HuggingFace checkpoint
\texttt{Qwen/Qwen2.5-7B}, quantized to FP8~E4M3 at engine startup with
per-tensor scale. The benchmark harness is \texttt{rvllm-v2-bench}
(the \texttt{crates/rvllm-v2/src/bin/bench.rs} binary), which drives
the engine directly without the HTTP layer, runs 3 iterations per
batch size after 1 warmup, and reports tokens per second as
$(\text{batch} \times \text{output\_len}) / \text{elapsed}$.
Greedy decoding (\texttt{temperature}=0, \texttt{ignore\_eos}=true),
512 output tokens per request, CUDA graphs enabled, all
\texttt{policy.json} entries pre-populated by
\texttt{autotune-cutlass}, \texttt{libfa3\_kernels.so} built from the
repository's build script. SHA: \texttt{eb9e247fd}. Date: April~16,
2026.

\subsection{Throughput by batch size}

\begin{tabular}{rrr}
\toprule
$\boldsymbol{N}$ & \textbf{tok/s} & \textbf{stdev} \\
\midrule
1   & 195.2     & 0.0 \\
4   & 770.2     & 0.7 \\
8   & 1{,}529.6 & 0.8 \\
16  & 3{,}011.8 & 0.6 \\
32  & 5{,}935.0 & 0.4 \\
64  & 11{,}064.3 & 1.6 \\
128 & \textbf{19{,}287.5} & 38.5 \\
\bottomrule
\end{tabular}

\subsection{Where the cycles go}

Per-kernel time (nsys profile, $N\!=\!32$, mid-decode), from an earlier
characterization of the same model on the same hardware:

\begin{itemize}
\item CUTLASS FP8 GEMM (fused residual, o\_proj): 15.9--65.6\,$\mu$s
  depending on bucket and variant.
\item FA3 SM90 paged decode: 7.5\,$\mu$s per layer at $N\!=\!32$.
\item fused RMSNorm + FP8 quant (per-token): 6.0--7.3\,$\mu$s.
\item fused SiLU*mul + FP8 quant: 14.3\,$\mu$s.
\item cuBLAS LM-head HGEMM: $\sim$483\,$\mu$s on $(32, 152064, 3584)$.
\end{itemize}

The LM-head HGEMM dominates a decode step at $N\!=\!32$. CUTLASS~FP8
alternatives for this shape were measured at $\sim$3\% slower than
cuBLAS at $M\!=\!128$ and are disabled in dispatch.

% ============================================================================
\section{Related work}
\label{sec:related}
% ============================================================================

vLLM~\cite{kwon2023vllm} introduced the paged-KV cache and continuous
batching patterns that rvLLM's scheduler implements. The paged KV
layout, block-manager structure, and preemption semantics follow vLLM's
conventions closely; the departure is the Rust runtime, the graph-replay
discipline, and the explicit rejection of the dispatch fallback chain.

FlashAttention-3~\cite{shah2024flash3} supplies the SM90 paged-decode
kernel; rvLLM links directly against a \texttt{.so} built from the
Hopper source tree and does not re-implement the attention kernel. The
\texttt{.so} build, variant selection, and workspace contract are
rvLLM's responsibility.

CUTLASS~\cite{cutlass2024} supplies the FP8 GEMM templates. rvLLM
generates the variant catalog, enforces the schedule-pairing rule, and
ships a per-shape autotune policy as a build artifact; the CUTLASS
library itself does not enforce pairings or manage workspaces across a
runtime.

TensorRT-LLM~\cite{tensorrt_llm} and SGLang~\cite{sglang2024} occupy the
same space but both assume a CUDA-Python host process. rvLLM's aim is
to demonstrate that a GPU-serving engine can be written entirely in a
systems language with no Python interpreter in the critical path, and
that doing so exposes the kernel-level correctness invariants that a
CUDA-Python host can paper over.

% ============================================================================
\section{Limitations and future work}
\label{sec:limitations}
% ============================================================================

The engine at \texttt{eb9e247fd} is single-GPU only. Multi-GPU tensor
parallelism exists in a sibling crate but is not wired into the v2 FP8
path. Pipeline parallelism is not implemented. The model catalog
validated end-to-end on this FP8 path is Qwen2 family, Llama 3 family,
Mistral 7B, and Gemma 2; other supported architectures compile but have
not been validated against HuggingFace reference on this version.
Speculative decode, LoRA serving, and tool-calling guided decoding are
present in the Python-free API layer but outside this paper.

The v3 rewrite consolidates the engine into 16 crates with a
DAG-enforced dependency graph and per-crate LoC budgets, and makes
each of the five invariants in Section~\ref{sec:invariants} either a
compile error, a type-level rejection, or a single-source-of-truth
runtime check. The CUDA \texttt{static\_assert} on schedule pairings,
the removal of the second metadata upload path, the
\texttt{GraphSafe} marker trait rejecting
\texttt{\&mut HbmArena} inside \texttt{CaptureScope} closures, and the
SHA-pinned \texttt{policy.json} artifact each eliminate one of the
three bug classes documented above. v3 is tracked under
\texttt{v3/SPEC.md} and \texttt{v3/IMPL\_PLAN.md} in the repository.

% ============================================================================
\section{Conclusion}
\label{sec:conclusion}
% ============================================================================

rvLLM demonstrates that an LLM inference engine competitive with modern
Python-hosted serving stacks can be written entirely in Rust, with a
single compute stream per worker, a frozen metadata layout per graph
bucket, and a kernel catalog in which every variant's workspace and
schedule are known at engine initialization. The engine achieves
19{,}287~tok/s serving Qwen2.5-7B in FP8 on a single H100, with
311~kernel launches per decode step compressed into one
\texttt{cuGraphLaunch} replay and a total of one DtoH and one HtoD
copy per step on the hot path.

The path to this number was not always direct: three failure modes
documented in Section~\ref{sec:invariants} produced either silent
correctness bugs or \texttt{CUDA\_ERROR\_ILLEGAL\_ADDRESS} only under
graph replay, and the engine spent one April day at roughly half the
peak throughput because of exactly these issues. The fixes, three
commits in one day, recovered the full performance envelope without
introducing fallbacks. The v3 rewrite makes each of the three bug
classes unrepresentable in the type system or caught at compile time.

% ============================================================================
\bibliographystyle{plainnat}
\bibliography{references}
% ============================================================================
